% TOP_PROBLEM: {"problemId": "problem.broues-abelian-defect-group-conjecture", "canonicalTitle": "Broues Abelian Defect Group Conjecture", "releaseRank": 337, "primaryDomain": "algebra_representation_category", "primaryDomainLabel": "Algebra, representation & category theory", "displayStatus": "open", "releaseStatus": "open"}
% !TeX program = lualatex
% Definition and sources only; this file is not an LLM solution attempt.
\documentclass[11pt]{article}
\usepackage[margin=25mm]{geometry}
\usepackage{fontspec}
\setmainfont{Latin Modern Roman}
\usepackage{amsmath,amssymb,mathtools}
\usepackage{unicode-math}
\setmathfont{Latin Modern Math}
\usepackage{xurl,hyperref}
\hypersetup{colorlinks=true,urlcolor=blue}
\setcounter{secnumdepth}{0}
\setlength{\parindent}{0pt}
\setlength{\parskip}{5pt}
\setlength{\emergencystretch}{3em}
\begin{document}
\section*{\texorpdfstring{Broues Abelian Defect Group Conjecture}{Broues Abelian Defect Group Conjecture}}\label{337}
\subsection{Definitions and mathematical statement}
Fix a splitting field $k$ of characteristic $p$ for a finite group $G$. A $p$-block algebra is $B=kGe$ for a primitive central idempotent $e$; its defect group $D$ is the associated maximal $p$-subgroup in Brauer's local block theory, defined up to conjugacy. Let $b$ be the Brauer correspondent of $B$ in $kN_G(D)$. If $D$ is abelian, Brou\'e's conjecture asks for a triangulated equivalence
\[
 D^b(B\text{-}\mathrm{mod})\simeq D^b(b\text{-}\mathrm{mod}).
\]
These bounded derived categories consist of bounded complexes of finitely generated modules, with quasi-isomorphisms inverted. The target is an equivalence of derived categories, not merely equality of character counts or a bijection between simple modules. Stronger refinements concerning a splendid equivalence are not imposed by the recorded statement.
\par\smallskip\textit{Formulation and concept references:} \href{http://www.numdam.org/item/AST_1990__181-182__61_0/}{[S1]}.

\subsection{Short English statement}
When a block has an abelian defect group, is its derived representation theory equivalent to that of its corresponding block in the defect group's normalizer?
\subsection{Sources}
\begin{itemize}
\item[S1]M. Broue, Asterisque 181-182:61-92, 1990.
\url{http://www.numdam.org/item/AST_1990__181-182__61_0/}.
\textit{Formulation source.}
\end{itemize}
\end{document}
